% ─────────────────────────────────────────────────────────────
% CHAPTER 5: DESIGN & IMPLEMENTATION OF BACKEND SERVICES
% ─────────────────────────────────────────────────────────────
\chapter{Design \& Implementation of Backend Services}
\label{ch:backend_implementation}

\section{Project Directory Structure and Modular Organization}
\label{sec:impl_structure}

To ensure long-term maintainability, testability, and adherence to clean code standards, the backend systems implemented at Digicon follow a strictly modular, layered directory hierarchy. Both the CRM Support Ticket platform and the Enterprise SMS Gateway microservice adhere to standardized architectural conventions. Listing~\ref{lst:project_structure} illustrates the standardized workspace layout.

\begin{lstlisting}[language=bash, caption={Modular Backend Project Directory Layout.}, label={lst:project_structure}]
src/
|-- common/                  # Cross-cutting concerns & shared utilities
|   |-- decorators/          # Custom parameter & metadata decorators (@Roles, @CurrentUser)
|   |-- filters/             # Global HTTP exception filters
|   |-- guards/              # Authentication & RBAC authorization guards
|   |-- interceptors/        # Logging, response transformation, & caching interceptors
|   |-- middlewares/         # Rate limiting, correlation IDs, & security headers
|-- config/                  # Environment variable validation & database configs
|-- modules/                 # Functional domain modules
|   |-- auth/                # Authentication, password hashing, & JWT issuance
|   |   |-- dto/             # LoginDto, RegisterDto, RefreshTokenDto
|   |   |-- auth.controller.ts
|   |   |-- auth.service.ts
|   |   |-- jwt.strategy.ts
|   |-- tickets/             # CRM Ticket lifecycle engine
|   |   |-- dto/             # CreateTicketDto, UpdateStatusDto, AssignTicketDto
|   |   |-- schemas/         # Mongoose Ticket & AuditLog document schemas
|   |   |-- ticket.controller.ts
|   |   |-- ticket.service.ts
|   |   |-- ticket.repository.ts
|   |-- sms/                 # Enterprise SMS gateway & queue dispatch
|   |   |-- dto/             # DispatchSmsDto, WebhookDlrDto
|   |   |-- entities/        # Prisma/TypeORM PostgreSQL entity mappings
|   |   |-- sms.controller.ts
|   |   |-- sms.service.ts
|   |   |-- sms.producer.ts  # BullMQ job enqueuer
|   |   |-- sms.worker.ts    # Background queue consumer
|   |-- cache/               # Redis in-memory cache-aside manager
|       |-- cache.service.ts
|-- main.ts                  # Application bootstrap, Swagger setup, & global pipes
\end{lstlisting}

\section{Implementation of Authentication, JWT, and RBAC Middleware}
\label{sec:impl_auth}

Stateless security is enforced across all API endpoints using JSON Web Tokens (JWT) coupled with declarative Role-Based Access Control (RBAC) guards.

\subsection{JWT Authentication Strategy}
When an incoming HTTP request reaches a protected route, the \texttt{JwtAuthGuard} intercepts the execution context and delegates token verification to the \texttt{JwtStrategy}. Listing~\ref{lst:jwt_strategy} presents the implementation of the passport-based JWT strategy.

\begin{lstlisting}[caption={Implementation of the Stateless JWT Verification Strategy.}, label={lst:jwt_strategy}]
import { Injectable, UnauthorizedException } from '@nestjs/common';
import { PassportStrategy } from '@nestjs/passport';
import { ExtractJwt, Strategy } from 'passport-jwt';
import { ConfigService } from '@nestjs/config';

export interface JwtPayload {
  sub: string;        // User UUID
  email: string;      // User Email Address
  role: string;       // Role identifier (e.g., 'TEAM_LEAD')
  permissions: string[];
}

@Injectable()
export class JwtStrategy extends PassportStrategy(Strategy) {
  constructor(private readonly configService: ConfigService) {
    super({
      jwtFromRequest: ExtractJwt.fromAuthHeaderAsBearerToken(),
      ignoreExpiration: false,
      secretOrKey: configService.get<string>('JWT_SECRET_KEY'),
    });
  }

  async validate(payload: JwtPayload) {
    if (!payload || !payload.sub) {
      throw new UnauthorizedException('Malformed or expired authentication token');
    }
    // Inject validated user context into Express Request object (req.user)
    return {
      userId: payload.sub,
      email: payload.email,
      role: payload.role,
      permissions: payload.permissions,
    };
  }
}
\end{lstlisting}

\subsection{Declarative RBAC Authorization Guard}
Endpoint access permissions are declared using custom TypeScript metadata decorators (e.g., \texttt{@Roles('SUPER\_ADMIN', 'TEAM\_LEAD')}). Listing~\ref{lst:roles_guard} demonstrates the execution guard that enforces this policy at runtime.

\begin{lstlisting}[caption={Declarative Role-Based Access Control (RBAC) Guard.}, label={lst:roles_guard}]
import { Injectable, CanActivate, ExecutionContext, ForbiddenException } from '@nestjs/common';
import { Reflector } from '@nestjs/core';

@Injectable()
export class RolesGuard implements CanActivate {
  constructor(private reflector: Reflector) {}

  canActivate(context: ExecutionContext): boolean {
    const requiredRoles = this.reflector.getAllAndOverride<string[]>('roles', [
      context.getHandler(),
      context.getClass(),
    ]);

    // If no specific roles are mandated, allow authenticated access
    if (!requiredRoles || requiredRoles.length === 0) {
      return true;
    }

    const { user } = context.switchToHttp().getRequest();
    if (!user || !user.role) {
      throw new ForbiddenException('Access denied: Unauthenticated user context');
    }

    const hasPermission = requiredRoles.includes(user.role);
    if (!hasPermission) {
      throw new ForbiddenException(
        `Insufficient privileges: Role '${user.role}' is unauthorized for this endpoint`
      );
    }

    return true;
  }
}
\end{lstlisting}

\section{Implementation of the CRM Support Ticket Routing Engine}
\label{sec:impl_crm}

The CRM Support Ticket platform handles customer issues arriving across Digicon's BPO contact center channels. The implementation leverages MongoDB via Mongoose to accommodate semi-structured customer case parameters.

\subsection{Ticket Controller Implementation}
The controller defines RESTful endpoints, enforces authorization guards, and validates client input through Data Transfer Objects (DTOs) decorated with \texttt{class-validator} annotations. Listing~\ref{lst:ticket_controller} details the controller methods for ticket ingestion and atomic agent assignment.

\begin{lstlisting}[caption={RESTful Ticket Management Controller.}, label={lst:ticket_controller}]
import { Controller, Post, Body, Patch, Param, UseGuards, Req, HttpStatus, HttpCode } from '@nestjs/common';
import { TicketService } from './ticket.service';
import { CreateTicketDto, AssignTicketDto } from './dto/ticket.dto';
import { JwtAuthGuard } from '../common/guards/jwt-auth.guard';
import { RolesGuard } from '../common/guards/roles.guard';
import { Roles } from '../common/decorators/roles.decorator';

@Controller('api/v1/tickets')
@UseGuards(JwtAuthGuard, RolesGuard)
export class TicketController {
  constructor(private readonly ticketService: TicketService) {}

  @Post()
  @HttpCode(HttpStatus.CREATED)
  @Roles('SUPER_ADMIN', 'TEAM_LEAD', 'SUPPORT_AGENT')
  async createTicket(@Body() createTicketDto: CreateTicketDto, @Req() req: any) {
    const creatorId = req.user.userId;
    return await this.ticketService.create(createTicketDto, creatorId);
  }

  @Patch(':id/assign')
  @HttpCode(HttpStatus.OK)
  @Roles('SUPER_ADMIN', 'TEAM_LEAD')
  async assignTicket(
    @Param('id') ticketId: string,
    @Body() assignTicketDto: AssignTicketDto,
    @Req() req: any
  ) {
    const supervisorId = req.user.userId;
    return await this.ticketService.assign(ticketId, assignTicketDto.agentId, supervisorId);
  }
}
\end{lstlisting}

\subsection{Ticket Service Business Logic and SLA Computation}
The service layer implements business rules, including automated SLA calculation and optimistic concurrency control. Listing~\ref{lst:ticket_service} illustrates the core service logic.

\begin{lstlisting}[caption={Ticket Service with SLA Calculation and Optimistic Locking.}, label={lst:ticket_service}]
import { Injectable, NotFoundException, ConflictException } from '@nestjs/common';
import { InjectModel } from '@nestjs/mongoose';
import { Model } from 'mongoose';
import { Ticket, TicketDocument } from './schemas/ticket.schema';
import { CreateTicketDto } from './dto/ticket.dto';
import { CacheService } from '../cache/cache.service';

@Injectable()
export class TicketService {
  constructor(
    @InjectModel(Ticket.name) private ticketModel: Model<TicketDocument>,
    private readonly cacheService: CacheService,
  ) {}

  async create(dto: CreateTicketDto, creatorId: string): Promise<Ticket> {
    const slaHours = this.calculateSlaWindow(dto.priority);
    const slaBreachTime = new Date(Date.now() + slaHours * 3600 * 1000);

    const ticketNumber = `TCK-${Date.now().toString().slice(-6)}-${Math.floor(100 + Math.random() * 900)}`;

    const newTicket = new this.ticketModel({
      ...dto,
      ticketNumber,
      status: 'OPEN',
      createdBy: creatorId,
      slaBreachAt: slaBreachTime,
      version: 1,
    });

    const savedTicket = await newTicket.save();
    
    // Invalidate cached ticket summaries
    await this.cacheService.invalidatePattern('tickets:summary:*');
    return savedTicket;
  }

  async assign(ticketId: string, agentId: string, supervisorId: string): Promise<Ticket> {
    const existing = await this.ticketModel.findById(ticketId);
    if (!existing) {
      throw new NotFoundException(`Ticket with ID ${ticketId} not found`);
    }

    // Optimistic Concurrency Control using version field
    const updated = await this.ticketModel.findOneAndUpdate(
      { _id: ticketId, version: existing.version },
      {
        $set: { assignedAgentId: agentId, status: 'IN_PROGRESS', updatedAt: new Date() },
        $inc: { version: 1 },
        $push: {
          auditLogs: {
            action: 'ASSIGNMENT',
            performedBy: supervisorId,
            assignedTo: agentId,
            timestamp: new Date(),
          },
        },
      },
      { new: true }
    );

    if (!updated) {
      throw new ConflictException('Concurrent update conflict: Ticket was modified by another user');
    }

    return updated;
  }

  private calculateSlaWindow(priority: string): number {
    switch (priority) {
      case 'CRITICAL': return 2;  // 2 hours SLA
      case 'HIGH':     return 6;  // 6 hours SLA
      case 'MEDIUM':   return 24; // 24 hours SLA
      default:         return 72; // 72 hours SLA for LOW
    }
  }
}
\end{lstlisting}

\section{Implementation of the Enterprise SMS Gateway Microservice}
\label{sec:impl_sms}

The SMS Gateway microservice is engineered to decouple high-volume outbound messaging from downstream carrier response latencies. It integrates the BullMQ message queue and a custom Token Bucket rate limiter.

\subsection{Token Bucket Rate Limiting Algorithm}
To adhere strictly to Bangladesh Telecommunication Regulatory Commission (BTRC) compliance and prevent carrier API flooding (HTTP 429), a Token Bucket rate limiter was implemented in Redis. The algorithm continuously replenishes tokens into a bucket at a fixed rate $R$ up to capacity $C$. Inbound dispatches consume tokens atomically via Lua scripts. Listing~\ref{lst:token_bucket} shows the implementation.

\begin{lstlisting}[caption={Token Bucket Rate Limiter Implementation in TypeScript.}, label={lst:token_bucket}]
import { Injectable } from '@nestjs/common';
import Redis from 'ioredis';

@Injectable()
export class TokenBucketLimiter {
  constructor(private readonly redis: Redis) {}

  async acquireToken(key: string, capacity: number, fillRatePerSec: number): Promise<boolean> {
    const now = Date.now();
    // Atomic Lua script ensures thread-safe token evaluation in Redis
    const luaScript = `
      local key = KEYS[1]
      local capacity = tonumber(ARGV[1])
      local fill_rate = tonumber(ARGV[2])
      local now = tonumber(ARGV[3])
      
      local data = redis.call('HMGET', key, 'tokens', 'last_updated')
      local tokens = tonumber(data[1])
      local last_updated = tonumber(data[2])
      
      if not tokens then
        tokens = capacity
        last_updated = now
      else
        local delta = math.max(0, (now - last_updated) / 1000)
        tokens = math.min(capacity, tokens + delta * fill_rate)
        last_updated = now
      end
      
      if tokens >= 1 then
        tokens = tokens - 1
        redis.call('HMSET', key, 'tokens', tokens, 'last_updated', last_updated)
        redis.call('EXPIRE', key, 60)
        return 1
      else
        redis.call('HMSET', key, 'tokens', tokens, 'last_updated', last_updated)
        return 0
      end
    `;

    const result = await this.redis.eval(luaScript, 1, `limiter:${key}`, capacity, fillRatePerSec, now);
    return result === 1;
  }
}
\end{lstlisting}

\subsection{BullMQ Queue Producer and Consumer Worker}
Outbound message requests are rapidly enqueued by the Producer service, as shown in Listing~\ref{lst:sms_queue}.

\begin{lstlisting}[caption={SMS Queue Producer and Asynchronous Worker Implementation.}, label={lst:sms_queue}]
import { Injectable, Logger } from '@nestjs/common';
import { Queue, Worker, Job } from 'bullmq';
import { ConfigService } from '@nestjs/config';
import axios from 'axios';

@Injectable()
export class SmsQueueProducer {
  private queue: Queue;

  constructor(private configService: ConfigService) {
    this.queue = new Queue('sms-dispatch-queue', {
      connection: {
        host: configService.get<string>('REDIS_HOST'),
        port: configService.get<number>('REDIS_PORT'),
      },
    });
  }

  async enqueueSms(recipient: string, message: string, trackingId: string) {
    await this.queue.add(
      'send-sms',
      { recipient, message, trackingId },
      {
        attempts: 3,
        backoff: {
          type: 'exponential',
          delay: 2000, // 2s, 4s, 8s backoff
        },
        removeOnComplete: true,
        removeOnFail: false,
      }
    );
  }
}

// Background Worker Service
@Injectable()
export class SmsQueueWorker {
  private worker: Worker;
  private readonly logger = new Logger(SmsQueueWorker.name);

  constructor(configService: ConfigService) {
    this.worker = new Worker(
      'sms-dispatch-queue',
      async (job: Job) => {
        const { recipient, message, trackingId } = job.data;
        this.logger.log(`Dispatching SMS to ${recipient} (Tracking ID: ${trackingId})`);
        
        // Execute HTTP POST to carrier SMPP/HTTP Gateway
        const response = await axios.post(configService.get('TELCO_GATEWAY_URL'), {
          msisdn: recipient,
          text: message,
          transaction_id: trackingId,
        }, { timeout: 4000 });

        return response.data;
      },
      {
        concurrency: 20, // 20 parallel worker threads per process
        connection: {
          host: configService.get('REDIS_HOST'),
          port: configService.get('REDIS_PORT'),
        },
      }
    );

    this.worker.on('failed', (job, err) => {
      this.logger.error(`Job ${job?.id} failed with error: ${err.message}`);
    });
  }
}
\end{lstlisting}

\subsection{Cryptographic Webhook Signature Verification}
When telecommunication operators send asynchronous Delivery Receipts (DLR), an HMAC-SHA256 signature middleware validates the payload authenticity to prevent spoofing, as demonstrated in Listing~\ref{lst:webhook_verifier}.

\begin{lstlisting}[caption={HMAC-SHA256 Webhook Cryptographic Verification Middleware.}, label={lst:webhook_verifier}]
import { Injectable, NestMiddleware, UnauthorizedException } from '@nestjs/common';
import { Request, Response, NextFunction } from 'express';
import * as crypto from 'crypto';

@Injectable()
export class WebhookSignatureMiddleware implements NestMiddleware {
  private readonly secret = process.env.TELCO_WEBHOOK_SECRET || 'secret-key';

  use(req: Request, res: Response, next: NextFunction) {
    const signature = req.headers['x-telco-signature'] as string;
    if (!signature) {
      throw new UnauthorizedException('Missing carrier cryptographic signature header');
    }

    const payload = JSON.stringify(req.body);
    const expectedSignature = crypto
      .createHmac('sha256', this.secret)
      .update(payload)
      .digest('hex');

    // Use timingSafeEqual to guard against timing attacks
    const isValid = crypto.timingSafeEqual(
      Buffer.from(signature, 'utf8'),
      Buffer.from(expectedSignature, 'utf8')
    );

    if (!isValid) {
      throw new UnauthorizedException('Invalid cryptographic signature: Tampered payload');
    }

    next();
  }
}
\end{lstlisting}

\section{Distributed Caching Implementation with Redis}
\label{sec:impl_caching}

To satisfy NFR-02 (sub-50ms latency), a generalized Cache-Aside service was implemented. The service manages serialized JSON data, TTL expiration, and pattern-based cache invalidation. Listing~\ref{lst:cache_service} presents the implementation.

\begin{lstlisting}[caption={Redis Cache-Aside Service Implementation.}, label={lst:cache_service}]
import { Injectable, Logger } from '@nestjs/common';
import Redis from 'ioredis';

@Injectable()
export class CacheService {
  private readonly redis: Redis;
  private readonly logger = new Logger(CacheService.name);

  constructor() {
    this.redis = new Redis({
      host: process.env.REDIS_HOST || '127.0.0.1',
      port: Number(process.env.REDIS_PORT) || 6379,
    });
  }

  async getOrSet<T>(key: string, ttlSeconds: number, fetcher: () => Promise<T>): Promise<T> {
    try {
      const cached = await this.redis.get(key);
      if (cached) {
        return JSON.parse(cached) as T;
      }
    } catch (err) {
      this.logger.warn(`Redis read error on key ${key}: ${err.message}`);
    }

    // Cache Miss: Execute primary database fetcher function
    const freshData = await fetcher();

    try {
      if (freshData !== null && freshData !== undefined) {
        await this.redis.set(key, JSON.stringify(freshData), 'EX', ttlSeconds);
      }
    } catch (err) {
      this.logger.warn(`Redis write error on key ${key}: ${err.message}`);
    }

    return freshData;
  }

  async invalidatePattern(pattern: string): Promise<void> {
    const keys = await this.redis.keys(pattern);
    if (keys.length > 0) {
      await this.redis.del(...keys);
      this.logger.log(`Invalidated ${keys.length} keys matching pattern: ${pattern}`);
    }
  }
}
\end{lstlisting}

\section{Global Error Handling and Centralized Logging}
\label{sec:impl_error_logging}

In enterprise environments, unhandled exceptions can leak internal stack traces and compromise security. A global HTTP exception filter was implemented to catch all anomalies and format them into standardized JSON error structures in conformance with RFC 7807 (Problem Details for HTTP APIs). Listing~\ref{lst:exception_filter} shows the global exception filter.

\begin{lstlisting}[caption={Global Centralized HTTP Exception Filter.}, label={lst:exception_filter}]
import { ExceptionFilter, Catch, ArgumentsHost, HttpException, HttpStatus, Logger } from '@nestjs/common';
import { Request, Response } from 'express';

@Catch()
export class GlobalExceptionFilter implements ExceptionFilter {
  private readonly logger = new Logger(GlobalExceptionFilter.name);

  catch(exception: unknown, host: ArgumentsHost) {
    const ctx = host.switchToHttp();
    const response = ctx.getResponse<Response>();
    const request = ctx.getRequest<Request>();

    const status =
      exception instanceof HttpException
        ? exception.getStatus()
        : HttpStatus.INTERNAL_SERVER_ERROR;

    const message =
      exception instanceof HttpException
        ? exception.getResponse()
        : 'Internal server error encountered';

    const errorPayload = {
      statusCode: status,
      timestamp: new Date().toISOString(),
      path: request.url,
      method: request.method,
      error: typeof message === 'object' ? message : { message },
    };

    this.logger.error(
      `HTTP ${status} [${request.method} ${request.url}]: ${JSON.stringify(errorPayload)}`
    );

    response.status(status).json(errorPayload);
  }
}
\end{lstlisting}
